\newcommand{\FinalMainComparisonCaption}{Performance at N=1024 on the original common500 cohort (500 sources, three noise realizations per source). Lines show published means; shaded bands are existing pointwise 95\% source-level bootstrap intervals. Failures remain included. Latent continuous JSCC uses the calibration-selected 100k checkpoint; adapted SwinJSCC uses the selected 200k checkpoint. Swin was initialized from 80k weights with fresh optimizer, sampling order and RNG, then trained for 120k additional updates over 1--19 dB. It stopped at budget while still improving; 19 dB is within this updated range. The old 80k adaptation had an out-of-range 19 dB point and is not plotted here. Adaptive-resolution BPG uses calibration-only selection. These later baseline comparisons reuse the original cohort, not a newly blind test. ConvNeXt agreement is agreement with the source-image prediction, not classification accuracy. DINOv2-L is ViT-L/14 cosine similarity. HiFi-DiffCom is omitted because no matching common500 result exists.}

\newcommand{\HiFiDevelopmentComparisonCaption}{Development-only qualitative examples at N = 1024 and SNR = 13 dB. Rows follow entries 5 through 8 of the pre-existing fixed16 source order (indices 4, 21, 24, 29). All columns match the same original uint8 source pixels. Neither source selection nor noise selection uses reconstruction quality. The fixed first noise realization is used: seed 2001 for the continuous, SwinJSCC and HiFi-DiffCom systems, and 6201 for digital transmission. SwinJSCC and HiFi-DiffCom share the exact received waveform; cross-family noise draws are different. These images are development examples and do not provide an independent evaluation on the unified 500-image holdout. All images come from verified existing reconstruction caches; no inference, channel simulation or metric evaluation was repeated. Columns show the original, adapted SwinJSCC image transmission, continuous latent-space joint source-channel coding, adapted HiFi-DiffCom diffusion-prior recovery, and partial-scale digital transmission with VAR completion. SwinJSCC uses the frozen 80k-step model and a 256-use protected header plus 768 body uses. HiFi-DiffCom uses this received frame and the frozen unconditional diffusion prior with its full posterior schedule. It has no label or text side information. The digital configuration at this point is m = 8, K = 9; its selected policy and frozen visual-model identity match the final release. No official-optimal or full-convergence claim is made for the adaptations.}

\newcommand{\HiFiDevelopmentMechanismCaption}{Development-only qualitative examples at N = 1024 and SNR = 13 dB. Rows follow entries 5 through 8 of the pre-existing fixed16 source order (indices 4, 21, 24, 29). All columns match the same original uint8 source pixels. Neither source selection nor noise selection uses reconstruction quality. The fixed first noise realization is used: seed 2001 for the continuous, SwinJSCC and HiFi-DiffCom systems, and 6201 for digital transmission. SwinJSCC and HiFi-DiffCom share the exact received waveform; cross-family noise draws are different. These images are development examples and do not provide an independent evaluation on the unified 500-image holdout. All images come from verified existing reconstruction caches; no inference, channel simulation or metric evaluation was repeated. Columns show the original, complete-scale transmission with VAR completion, partial-scale transmission without generative completion, and partial-scale transmission with VAR completion. Complete-scale means the selected prefix of complete scales (m = 8, K = 0, 255 tokens), not all source scales. Partial-scale transmission sends m = 8, K = 9 (264 tokens). Both digital configurations have the same modulation and coding permissions. The last two columns use the same actual received tokens and frozen decoder Dc; absent residual contributions are zero in direct decoding, not codebook index zero, and no extra source-image ground truth is supplied.}
